%!TEX root = ../../../main.tex
%
% §8.4.1 H-01 診断の再掲と FCCD 救済（CHK-193 Phase D-1 で本編から退避）
%
% 退避理由:
%   §4.5 04e_fccd.tex の H-01 一次定義（sec:fccd_motivation_h01）と内容完全重複．
%   E-030 Exp-1 の $|\BF_\text{res}| \approx 884$ at step 1 / $\rho_l/\rho_g=830$ の
%   数値まで両所で独立に再掲していた．
%
% 本編の代替:
%   paper/sections/08_2_fccd_bf.tex L19-40 → 3 行の橋渡し段落 + cross-ref
%   `§\ref{sec:fccd_motivation_h01}` に圧縮．
%   外部参照（10_5_pure_fccd_dns, 08c_pressure_filter, 10_3_level_selection,
%   appendix_gpu_fvm）のため sec:h01_diagnosis_fccd_remedy は phantomsection 化．
%
% 再活性化条件:
%   H-01 診断を §4.5 と独立の議論として §8.4 で完全展開する必要が生じた場合．
%   本編から再 \input するか，内容を L18 付近に復元する．

% -------------------------------------------------------------------------------
\subsubsection{H-01 診断の再掲と FCCD 救済}
\label{sec:h01_diagnosis_fccd_remedy}
% -------------------------------------------------------------------------------

非一様格子上で節点 CCD 勾配と FVM の面平均勾配 $\mathcal{G}^\text{adj}$
は異なる離散幾何を参照する：
\begin{equation}
  \mathcal{G}^\text{adj} p\big|_{i+1/2}
    = \frac{p_{i+1}-p_i}{d_f^{(i+1/2)}},
  \qquad d_f^{(i+1/2)} = \text{局所面距離．}
  \label{eq:bf_gadj_recap}
\end{equation}
$\mathcal{G}^\text{adj}$ は速度投影で保存性を保証する一方，
節点 CCD 勾配は $\sigma\kappa\bnabla\psi$ に使われ，
\begin{equation}
  \mathrm{BF}_\text{res}
    := \bigl\| \mathcal{G}^\text{adj} p_\text{YL} - \sigma\kappa\bnabla\psi\,\text{（CCD）} \bigr\|_\infty
    \sim \Ord{\Delta x^2}\,\Bigl|\frac{d\log J}{dx}\Bigr|
  \label{eq:bf_residual_h01}
\end{equation}
が計測される（Exp-1：$\rho_l/\rho_g{=}830$ で $|\BF_\text{res}|\approx 884$ at step 1）．
$\sigma{=}0$（Exp-2）で爆発が消える事実がH-01 を確定させた．

\noindent\textbf{FCCD 統一}：
FCCD 演算子 $D^\FCCD$ は $f=i-1/2$ で $\{u_{i-1},u_i\}$ のみを用い，
$\mathcal{G}^\text{adj}$ と\textbf{同一の面位置・同一のデータフットプリント}．
\begin{equation}
  \mathcal{G}p := D^\FCCD p,
  \qquad \bnabla\psi := D^\FCCD\psi
  \label{eq:bf_fccd_unification}
\end{equation}
を採れば，BF 残差は FCCD 切断誤差に還元され
\begin{equation}
  \mathrm{BF}_\text{res} = \Ord{\Delta x^4} \ \gg\ \text{従来 }\Ord{\Delta x^2}.
  \label{eq:bf_residual_fccd}
\end{equation}
